\chapter{バックテストと過学習}
\label{ch:backtest}

\textbf{バックテスト}とは、戦略を過去のデータに当てはめ、もしその戦略で運用していたらどうなったかを計算することである。新しい戦略の検証のほとんどは、バックテストに頼るしかない。しかし、バックテストの成績は、ほぼ必ず、実際の運用の成績より良く出る。この章は、その理由と対策をまとめる。

\section{バックテストの偏り}

\begin{table}[H]
\centering\small
\caption{バックテストの成績を良く見せる偏り}
\label{tab:biases}
\begin{tabular}{@{}P{0.22\linewidth}P{0.42\linewidth}P{0.28\linewidth}@{}}
\toprule
偏り & 中身 & 対策 \\
\midrule
先読み & その時点で知られていなかった情報を使う（決算の公表日の前に決算の値を使う、後で訂正された値を使う、言語モデルが学習した後の出来事を知っている） & ポイント・イン・タイムのデータ、公表の遅れ、学習の期間の止まった言語モデル \\
生存者 & 今日存在する銘柄・ファンドだけで過去を検証する & 上場廃止・倒産した銘柄を含むデータ \\
データ・スヌーピング & 多くの戦略・パラメータを試し、最も良いものを報告する & 試した数を記録し、多重検定の補正（\ref{sec:mt} 節） \\
コストの無視 & 取引コスト・空売りの費用・インパクトを入れない & \chapref{ch:cost} のコストのモデル \\
実行できない売買 & 終値で判断し同じ終値で売買する、売買の止まった銘柄を売買する、借りられない株を空売りする & 判断と売買の時点をずらす、流動性の条件 \\
期間の選び方 & 戦略に都合の良い期間だけを使う & できるだけ長い期間、複数の市場 \\
\bottomrule
\end{tabular}
\end{table}

\section{過学習}

\subsection{多くを試すと、偶然に良いものが見つかる}

本当は予測力がまったくない（真のシャープレシオがゼロの）戦略を $N$ 個試すと、その中の最大のシャープレシオは、$N$ とともに大きくなる。Bailey・López de Prado（2014）\cite{dsr2014} によれば、試した戦略のシャープレシオの推定値のばらつきを $\sigma_{\SR}$ とすると、最大値の期待値はおよそ
\begin{equation}
\E[\max_n \widehat{\SR}_n] \approx \sigma_{\SR}\Big[(1-\gamma)\,\Phi^{-1}\!\Big(1 - \frac{1}{N}\Big) + \gamma\,\Phi^{-1}\!\Big(1 - \frac{1}{N e}\Big)\Big]
\end{equation}
である（$\gamma \approx 0.5772$ はオイラーの定数）。

\begin{exbox}[偶然でどれだけ良く見えるか]
5 年の日次のデータで、真のシャープレシオがゼロの戦略を試す。年率のシャープレシオの推定値の標準誤差は約 $1/\sqrt{5} \approx 0.45$。
\begin{center}
\begin{tabular}{@{}lcccc@{}}
\toprule
試した数 $N$ & 10 & 100 & 1{,}000 & 10{,}000 \\
\midrule
最大の SR の期待値（年率） & 0.70 & 1.13 & 1.46 & 1.73 \\
\bottomrule
\end{tabular}
\end{center}
100 通りのパラメータを試せば、まったく無意味な戦略から、年率 1 を超えるシャープレシオが「見つかる」。Bailey・Borwein・López de Prado・Zhu（2014）\cite{charlatanism2014} は、試した数を示さずにバックテストの成績を示すことを、「疑似数学と金融のいかさま」と呼んで警告した。
\end{exbox}

\subsection{デフレートしたシャープレシオ}

Bailey・López de Prado（2014）の\textbf{デフレートしたシャープレシオ}（DSR）は、観測したシャープレシオ $\widehat\SR$ が、「$N$ 回試したときに偶然出る最大値」$\SR_0$ を超える確率を、リターンの歪度 $\gamma_3$・尖度 $\gamma_4$ と標本の長さ $T$ も考慮して計算する。
\begin{equation}
\mathrm{DSR} = \Phi\!\left(\frac{(\widehat\SR - \SR_0)\sqrt{T-1}}{\sqrt{1 - \gamma_3\,\widehat\SR + \frac{\gamma_4 - 1}{4}\widehat\SR^2}}\right)
\end{equation}
（$\widehat\SR$ と $\SR_0$ は期間あたりの値。$\SR_0$ は上の最大値の期待値。）これが 0.95 を超えなければ、偶然と区別できない。重要なのは、\textbf{試した数 $N$ を記録しておかなければ、この補正ができない}ことである。

\subsection{バックテストの過学習の確率}

Bailey・Borwein・López de Prado・Zhu（2017）\cite{pbo2017} の\textbf{PBO}（probability of backtest overfitting）は、次の手順で計算する（組み合わせ対称の交差検証、CSCV）。

\begin{enumerate}
\item 期間を $S$ 個（例えば 16）の塊に分ける。
\item そのうち半分を学習、残りの半分を評価とする分け方を、すべての組み合わせ（$\binom{16}{8} = 12{,}870$ 通り）で作る。
\item 各分け方で、試したすべての戦略のうち、学習の期間で最も良いものを選び、それが評価の期間で何番目かを見る。
\item 評価の期間で中央値より下になった割合を PBO とする。
\end{enumerate}
PBO が 0.5 に近いなら、「学習の期間で最も良い」ことに、ほとんど意味がない。

\section{多重検定}
\label{sec:mt}

\subsection{補正の方法}

$m$ 個の仮説を同時に検定するとき、それぞれを 5\% で検定すると、少なくとも一つが偶然に有意になる確率は $1 - 0.95^m$（$m = 20$ なら 64\%）になる。

\begin{itemize}
\item \textbf{Bonferroni}：各検定を $\alpha/m$ で行う。最も厳しい。
\item \textbf{Holm}\cite{holm1979}：$p$ 値を小さい順に並べ、$k$ 番目を $\alpha/(m-k+1)$ と比べ、初めて超えたところで止める。Bonferroni より検出力が高く、同じく「一つでも誤る確率」（FWER）を $\alpha$ 以下に抑える。
\item \textbf{Benjamini--Hochberg}\cite{bh1995}：有意とした仮説のうち、誤りの割合の期待値（FDR）を抑える。多くの候補から有望なものを拾うときに向く。
\item \textbf{White の Reality Check}（White 2000 \cite{white2000}）、\textbf{Hansen の SPA}（Hansen 2005 \cite{hansen2005}）：多くの戦略の中で最も良いものが、基準（例：市場）に本当に勝っているかを、ブートストラップで検定する。
\end{itemize}

\subsection{\texorpdfstring{$t > 3$}{t > 3}}

Harvey・Liu・Zhu（2016）\cite{hlz2016} は、学術誌に発表された数百のファクターについて、研究者の全体が試した数を考慮すると、新しいファクターの $t$ 値の基準は 2 ではなく 3 程度にすべきだと論じた。業界でも、新しいシグナルの採用の基準として広く引用される。

\section{検証の設計}

\subsection{手元に残しておく}

最も強い対策は、\textbf{一度も見ていないデータ}で、\textbf{一度だけ}評価することである。

\begin{enumerate}
\item データを、開発の期間と、最後まで見ない期間（ホールドアウト）に分ける。
\item 開発の期間の中で、ウォークフォワードと交差検証（\ref{sec:cv} 節）で戦略を作る。
\item 戦略を固定してから、ホールドアウトで一度だけ評価する。ホールドアウトの結果を見て戦略を直したら、それはもうホールドアウトではない。
\end{enumerate}

\subsection{事前登録}

医学の臨床試験と同じく、検定の仮説・データ・変数・方法・判定の基準を、結果を見る前に文書にして固定する（事前登録）。事前登録にない分析は「探索」として区別する。分析の選択の自由（どの期間、どの統制、どの外れ値の処理）による偶然の発見を防ぐ。

\subsection{紙の上の運用と小さな実運用}

バックテストの後、実際の価格で売買の注文を（出さずに）記録する「紙の上の運用」や、小さな資金での実運用を、数か月〜数年行う。実際の売買のコスト、データの遅れ、システムの誤りは、ここで初めて分かる。

\section{バックテストの確認の一覧}

\begin{keybox}[バックテストを信じる前に確かめること]
\begin{enumerate}
\item 各時点で、その時点で知られていた情報だけを使っているか（公表日、訂正、言語モデルの学習の期間）。
\item 上場廃止・倒産した銘柄が含まれているか。
\item 取引コスト・空売りの費用・インパクトを入れたか。容量はどれくらいか。
\item 何通りの戦略・パラメータを試したか。それを補正したか。
\item 期間を分けても（前半と後半、好況と不況）、市場を変えても、同じ結果か。
\item 結果に経済的な説明があるか（なぜ儲かるのか、誰が損をしているのか）。
\item 少数の銘柄・少数の日が、収益の大半を生んでいないか。
\item 最も良いパラメータの周りで、成績がなだらかに変わるか（鋭い山なら過学習）。
\end{enumerate}
\end{keybox}

\begin{projbox}
LLM-BLM の小型株の実験 S1 は、検定の仮説・変数・判定の基準・コードを、株価の反応を一度も計算する前にコミットして固定し（事前登録）、1 回だけ実行する設計にした。副次の検定には Holm の補正をかけ、事前登録にない分析は「探索」として別に報告する。LLM の点数は、学習の期間が止まった言語モデル（ChronoGPT の 2015 年版）を主に使い、言語モデルが将来を知っている先読みを避けた。これらは、この章の対策をそのまま当てはめたものである。
\end{projbox}
